%% ===================================================================
\section{Numerical checks}\label{app:checks}
%% ===================================================================

The analytic results were checked numerically in two independent ways.

\paragraph*{Operator model of the leading logarithms.} We built an exact operator representation of Eq.~\eqref{eq:xsec} at LL. The valence charges are two or three quarks and antiquarks, with $3\times3$ colour matrices. The Wilson lines are random $SU(3)$ matrices at all transverse points. The measured and the soft gluon live in explicit Fock spaces with two polarizations, and the emission operators~\eqref{eq:OmegaLLfact} are expanded to the required order. Table~\ref{tab:checksLL} lists the comparisons.

\begin{table}[!htbp]
\centering
\small
\renewcommand{\arraystretch}{1.2}
\begin{tabular}{p{9.6cm}p{5.4cm}}
\toprule
check & result\\
\midrule
LO: operator calculation vs Eq.~\eqref{eq:Oop} & ratio $1$\\
Region T: REAL$+$VIRTUAL vs the double commutator~\eqref{eq:HLR} & residual $6\times10^{-14}$\\
Region T: explicit form of App.~\ref{app:explicit} vs operator calculation & agreement to machine precision\\
Region T: $\sum$ of the LL rows of Group~I vs exact virtual part; $\sum$ of Groups~II$+$III vs exact real part & both $10^{-13}$\\
Region T: every row of App.~\ref{app:LLrows}, symbolic vs numerical & $10^{-14}$ (18 rows)\\
Region T: surviving two-charge classes vs Eq.~\eqref{eq:HLR}, per kernel & agreement for all 8 kernels\\
Region P: operator calculation vs Eq.~\eqref{eq:LOBFKL} (colour-singlet $q\bar q$) & ratio $1.0000000$\\
Region P: $\sum$ of rows vs (i)$+$(ii)$+$(iii) (charges $q,q,\bar q$) & $4\times10^{-14}$\\
Region P: (i)$+$(ii)$+$(iii) vs exact operator result & $9\times10^{-14}$\\
colour identities of Sec.~\ref{sec:T} (Step 6) & $10^{-16}$\\
\bottomrule
\end{tabular}
\caption{Checks of the leading-log results against an exact operator model.}
\label{tab:checksLL}
\end{table}

The same model shows that two expressions of the draft notes require the corrections incorporated here. In the explicit JIMWLK$\otimes$LO, the matrix $M$ must be evaluated at the points where the derivatives act, and $dN^{(5)}$ enters with a positive sign; without these corrections the sum differs from $-H_{\rm JIMWLK}\otimes$LO by a factor $2.13$ at a test configuration. In addition, the three-charge LL terms of G3-III part~1 do not reproduce the exact real term (Sec.~\ref{sec:summary}).

\paragraph*{Longitudinal integrals and limits.} For random transverse configurations we checked the following.
\begin{itemize}
\item The analytic $p^+$ integrals of G1-III, of the $\mf B_2$ structure $\mathfrak F$ and of G2-I(2) agree with direct numerical quadrature, with the principal value at the pole, to 9--11 significant digits.
\item The split~\eqref{eq:findef} is exact: finite part plus subtracted logarithms reproduces the full row to 10--12 digits at several $(\Lambda,\vee)$.
\item G1-III with the sign~\eqref{eq:nonInst1}: no term grows like $1/\Lambda$. The full integral, evaluated with 30-digit arithmetic at $\vee=10^3$ and $\Lambda=10^{-3},10^{-5},10^{-7}$, grows like $\log(1/\Lambda)$ with slope $0.1076475$, compared with $(K_b-2K_f)/k^+=0.1076479$ predicted by Eq.~\eqref{eq:LLIII}. The finite part~\eqref{eq:finIII} converges to $-0.78197$ (Sec.~\ref{sec:finiteG1}).
\item G1-II(2): the coefficient of $\log\vee$ vanishes to $10^{-10}$. G2-I(2): the $\delta Y_P$ coefficients of the $\KK_A$, $\KK_B$, $\KK_D$, $\KK_E$ and $d_\perp$ terms vanish identically (computer algebra).
\item Eq.~\eqref{eq:phase} and the Fourier transforms~\eqref{eq:i3}--\eqref{eq:i5}, by quadrature.
\item The collinear and UV limits of Table~\ref{tab:dglap}, at $d_\perp=2$ and $d_\perp=1.6$. G2-I(2) gives $\Phat(\xi)$ to six digits and G1-V(2b) gives $\Phat(\xi)$. G3-III(2) gives $-1$ times the G2-I(2) limit, both from the integrand before and after the $w'$ integration. G2-IV gives $+1$ times the G2-I(2) limit. The self-energy integrand equals $\Phat$, and Eq.~\eqref{eq:intPhat} holds (computer algebra).
\item With a general adjoint $U$ and the index orders of the draft notes, the colour operators of G2-I(2), G2-IV and G3-III(2) at $w=w'=z$ differ (Sec.~\ref{sec:dglap}).
\end{itemize}
